\documentclass[11pt]{article}

\usepackage[utf8]{inputenc}
\usepackage[T1]{fontenc}
\usepackage[margin=1in]{geometry}
\usepackage{xcolor}
\usepackage{hyperref}
\usepackage{helvet}
\usepackage{titlesec}
\usepackage{microtype}

\renewcommand{\familydefault}{\sfdefault}
\pagestyle{empty}

\titleformat{\paragraph}[runin]{\bfseries}{}{}{}[.]
\titlespacing{\paragraph}{0pt}{8pt}{6pt}

\setlength{\parskip}{6pt}
\setlength{\parindent}{0pt}

\begin{document}

\begin{center}
{\Large\bfseries Research Statement}\\[4pt]
{\bfseries Common-Noise Mean-Field Games: Theory, Computation, and Financial Applications}\\[4pt]
Nisong Monyimba \;|\; nmonyimb@asu.edu
\end{center}
\vspace{8pt}

My research investigates mean-field games (MFGs) in the presence of common noise --- aggregate
shocks that affect every agent in a population simultaneously, as opposed to the idiosyncratic,
individually-averaged noise of the classical MFG setting. This single structural change transforms
the population's distribution from a deterministic object into a stochastic process in its own right,
and it is the source of nearly every technical difficulty I have spent the past year working through:
in well-posedness theory, in numerical schemes, in reinforcement-learning algorithms that attempt to
learn equilibria from data, and in financial applications where the population's distributional state
is itself the object of economic interest.

\paragraph{How I arrived here} My background is in biomedical engineering --- a B.Sc.\ from KNUST
and an M.Sc.\ from Arizona State University, followed by several years in industry as a quality and
process engineer (Abbott Laboratories, Vastek, and earlier at Mayo Clinic/ASU), where I worked
extensively with statistical process control, Monte Carlo simulation, design of experiments, and
risk modeling. That work built real fluency with applied probability and numerical methods, but it
did not include the formal, proof-based mathematical training --- real analysis, measure theory,
rigorous stochastic processes --- that underlies the literature I have since become genuinely
absorbed in. Over the past year, working independently and largely outside any formal coursework,
I built a complete master's-thesis-length research program in common-noise mean-field games: a
158-page manuscript spanning quantitative well-posedness theory, numerical analysis, multi-timescale
reinforcement learning, and two financial applications (systemic risk with optimal government
bailouts, and a Wasserstein-distance-based measure of market ``crowding'' validated empirically).
I am applying to PhD programs because I want --- and currently lack --- the rigorous theoretical
foundation to carry this work forward independently, defend it at the level it deserves, and extend
it into open questions I can currently only describe, not yet resolve myself.

\paragraph{What I have done, stated precisely} The thesis's central technical contribution is an
explicit, parameter-dependent contraction estimate for the mean-field fixed-point map governing
common-noise equilibria, obtained via a synchronous-coupling argument that isolates how the
common-noise coefficient enters the stability constant through a second-order BSDE effect rather
than a first-order state-difference term. Building on this, I developed a combined
particle/finite-difference numerical scheme with a proven convergence rate and an explicit
CFL-type stability condition for the coupled system; a three-timescale reinforcement-learning
algorithm with a subsequential convergence guarantee, together with a hierarchical actor-critic
architecture motivated by an SPDE decomposition of the associated master equation; a closed-form
Stackelberg model of systemic risk and optimal government bailouts; and a novel ``crowding
measure'' --- the Wasserstein-2 distance between the observed and theoretical equilibrium
distributions --- as a candidate predictor of asset returns.

\paragraph{What I want to be equally precise about} I validate this work computationally and
through careful, repeated empirical checks --- including, over the course of this project, catching
and correcting a genuine sign error in a closed-form Riccati system that had silently produced
non-optimal (though stable and plausible-looking) numerical results across several downstream
experiments, a correction I only found by insisting on an independent optimality check rather than
trusting that a well-behaved simulation was a correct one. That experience shaped how I think about
this work: I am confident in my ability to build, test, and rigorously interrogate computational
models of stochastic systems, and to recognize the difference between code that runs cleanly and a
result that is actually correct. I am not yet able to independently derive or defend every proof in
this thesis from first principles without reference material, and I do not want to overstate that I
am. This is precisely the gap I am seeking a PhD to close: I have begun working independently
through real analysis and rigorous stochastic processes in parallel with this application cycle, and
I view the structured theoretical training of a PhD program --- not just its research opportunities
--- as essential, not incidental, to where I want this research program to go next.

\paragraph{Where I want to take this} Several concrete open directions emerged directly from this
work: extending the common-noise framework to jump-diffusion shocks, where the
synchronous-coupling argument I rely on does not directly apply; relaxing the strong isoperimetric
(log-Sobolev) condition needed for the reinforcement-learning convergence rate; replacing the
heuristic algorithm-to-operator mapping in the hierarchical RL architecture with a more principled
correspondence; and --- most concretely --- recalibrating the empirical alpha-generation results
against real institutional-holdings and market data, which I was unable to access in the
environment where this thesis was built. I am also interested in the formal verification direction
I began in Lean 4 for a handful of the thesis's key estimates, though I want to be equally honest
that this was exploratory and not machine-checked against a live proof assistant in this iteration
of the work.

\paragraph{Fit} My thesis's numerical contribution --- a combined particle/finite-difference
scheme for the coupled FBSDE-SPDE system underlying common-noise mean-field game equilibria ---
overlaps directly with Professor Zhang's research in numerical approximation of stochastic games
and optimal stochastic control. I am particularly interested in his work on adaptive finite-element
approaches for related nonlinear PDEs as a possible route to relaxing the smoothness assumptions
my current convergence theorem requires.

\end{document}
